% paper/tables/v3_design.tex -- generated by analyze/paper/tables_paper.py; sources in paper/tables/README.md
\setlength{\tabcolsep}{3pt}\renewcommand{\arraystretch}{0.85}%
\begin{tabular}{@{}p{0.16\linewidth}p{0.4\linewidth}llp{0.25\linewidth}@{}}
\toprule
条件・データ & 中身 & 組 & 実験 & 位置づけ \\
\midrule
経験なし・50\%・80\% & A-1 が規則の一歩を選ぶ確率 $p=0$，0.5，0.8 で段階Iを学習 & 組1〜6 & 1〜3 & 事前登録（判定は組1〜3） \\
符号無作為（経験80\%） & 経験80\%と同じ軌跡で，$P_t$ の符号だけをエピソードごとに無作為に反転 & 組1〜6 & 1・2 & 組1〜3 は事前登録，6組の判定は組1〜3 の後に登録 \\
符号無作為（経験50\%） & 経験50\%と同じ軌跡で，同じエピソードの $P_t$ の符号を反転 & 組1〜6 & 1・2 & 6組の判定を学習の前に登録 \\
非共有条件 & LSTM を process-1 用と process-2 用に分ける（経験80\%のデータ） & 組1〜3 & 1 & 事前登録した対照 \\
零入力条件 & 経験80\%と同じ軌跡で $P_t$ を常に0 & 組1〜3 & 1 & 補助の対照（入口の重みは学習されない） \\
好み切替条件 & 好きな色をエピソードごとに赤か緑にし，$P_t$ ＝ 好きな色 $-$ 嫌いな色（$p=0.5$） & 組1〜3 & 1 & 補助．$d_P$ は事前登録外の確認 \\
$G$（自由出力） & 段階Iの予測器（経験なし・経験50\%）の上で，好み混合データで予測誤差だけから学習 & 組1〜4 & 3 & 段階Cの主の対象（T5 と同じ $G$） \\
学習前の $G$ & 段階IIの学習を始める直前に保存した $G$（経験50\%の予測器の上） & 組1〜4 & 3 & 事前登録（Q1c．第1段の不成立により記述） \\
\multicolumn{2}{@{}p{0.58\linewidth}}{自分の学習データと同じ作りの評価データ（A-1 は規則の確率 0.8）} & 組1〜6 & 1 & 計算の前に判定の基準を記録 \\
\multicolumn{2}{@{}p{0.58\linewidth}}{段階Bの評価データ（新しい鍵，配置×好みの4セルに 300 ずつ，計 1200）} & ― & 2・3 & データの確認とともに作る前に登録 \\
\midrule
指標 & \multicolumn{2}{l}{何を測るか} & 実験 & 事前登録の扱い \\
\midrule
$g$（G3） & \multicolumn{2}{p{0.475\linewidth}}{入力を隠し快・不快を保つと，予測が赤の方へ進み続ける傾き} & 1 & 事前登録した主 \\
$d_P$ & \multicolumn{2}{p{0.475\linewidth}}{$(g_{\mbox{\scriptsize 保つ}}-g_{\mbox{\scriptsize 反転}})/2$．自分側の快・不快への依存} & 1 & 事前登録（符号無作為との比較） \\
G5b & \multicolumn{2}{p{0.475\linewidth}}{process-2 の一範囲への $\pm a$ による，予測の中の A-2 のずれの差} & 1 & 事前登録した主 \\
符号に従う分 & \multicolumn{2}{p{0.475\linewidth}}{本当の快・不快を保った傾きと，符号を反転した傾きの差の半分} & 1 & 定義は事後，組5・6 と符号無作為の比較は事前登録 \\
$U_{\mathrm{pref}}$，$U^{\mathrm{pos}}_{\mathrm{pref}}$ & \multicolumn{2}{p{0.475\linewidth}}{他者自身の視界から作る候補で，反対の好みと正しい好みの他者側の誤差の差（$U^{\mathrm{pos}}_{\mathrm{pref}}$ は予測画像の A-2 の位置のずれの差．診断）} & 2 & 事前登録 \\
$b_{\mathrm{pref}}$，$b_{\mathrm{dir}}$（段階B） & \multicolumn{2}{p{0.475\linewidth}}{観察者が作れる候補 $c^{\pm}$ での誤差の差 $Q_e$ の，好みの側・移動の側の係数（赤の側・初期位置をそろえる）} & 2 & 事前登録（族 P は主，族 D は副．(v) の経験の比較は事前登録外） \\
T2・T5 & \multicolumn{2}{p{0.475\linewidth}}{他者用モジュールの出力 $Q2$ が好みや色に応じて変わるか} & 3 & 事前登録した主 \\
$b^{q}_{\mathrm{pref}}$（段階C (L1)） & \multicolumn{2}{p{0.475\linewidth}}{$Q2$ を好みの軸 $c^{+}$ に射影した値（標準化）の好みの係数} & 3 & 事前登録した主（第1段：Q1a・Q1b） \\
読み出し（段階C (L2)・(L3)） & \multicolumn{2}{p{0.475\linewidth}}{$Q2$（L2）と，$G$ の入力・内部状態・$h^2$（L3）から好みを読む BA・$\Delta\mathrm{LL}_{\mathrm{pref}}$・一般化} & 3 & (L2) は副，(L3) は記述．正則化の選び方は実データの集計の前に修正．一般化の差は事前登録外 \\
\bottomrule
\end{tabular}
